% Section: Discussion
% Source: paper_draft_v0.1.md § 6. Discussion
% [VERIFY ACM TEMPLATE BEFORE SUBMISSION]

\section{Discussion}
\label{sec:discussion}

\subsection{Interpretation of the Negative-Result Finding}

The 86.2\% failure rate is not a finding against ML in security operations broadly.
It is a finding about the relative performance of unsupervised ML versus structured
rule baselines on \emph{synthetic, structured, class-imbalanced tabular anomaly
detection} in the specific domain of cyber-hygiene telemetry.

The synthetic benchmark is, by design, favorable to rule-based methods: anomaly
injection creates features that are directly correlated with the anomaly class in
ways that rule thresholds exploit cleanly. In operational environments, data is
noisier, labels are unavailable, and the feature space is more complex.
\hygienebench reports this limitation explicitly in the dataset card and generator
documentation.

The appropriate takeaway is: for \hygienebench v0.1, ML adds consistent value over
rule baselines specifically on T2 (group membership drift) and T5
(patch/vulnerability hygiene). These tasks involve multi-signal combinations that
benefit from learned representations. Future work should test whether this advantage
persists under increased population-level noise and label sparsity.

\subsection{Limitations}

\paragraph{Synthetic data.}
All telemetry is synthetic. Structural priors are approximate (DBIR 2026 for patch
lag; NVD for CVE severity; CISA BOD 23-01 for freshness cadence). AD group-size
distributions lack a peer-reviewed source and are modeled under disclosed
assumptions. Detection performance on real operational data will differ from
benchmark results.

\paragraph{Approximations in M6 and M8.}
M6 is implemented as PCA-based linear reconstruction error (not a neural
autoencoder) and M8 as networkx graph features + Isolation Forest (not
DOMINANT-style deep graph autoencoder~\cite{ding2019dominant}). These are disclosed
approximations; results may change with full implementations.

\paragraph{Small population scale.}
Medium scale ($n$=1,000 users) limits the statistical power of the evaluation. The
population is smaller than most public-sector environments, which may affect the
generalizability of the class imbalance findings.

\paragraph{Feature engineering as a design choice.}
The task feature extractors (T1--T7) encode domain knowledge about which features
matter. Methods that exploit these features more effectively will appear better. A
fully unsupervised method that discovers relevant features from raw tables is not
evaluated.

\paragraph{Fixed $k$ values.}
P@$k$ is sensitive to the choice of $k$. \hygienebench uses operationally motivated
$k$ values (daily/weekly SOC review budgets) rather than tuning $k$ to maximize
method performance. This is a deliberate design choice but limits comparison to
benchmarks with different $k$ conventions.

\paragraph{No calibration.}
Anomaly scores are not calibrated to probabilities. Calibration is the primary
research question of a separate line of work and is excluded from \hygienebench by
design.

\subsection{Relationship to Prior Work}

\hygienebench is designed to complement, not replace, existing benchmarks. It does
not claim to detect real network intrusions (LANL, CICIDS), insider threats (CERT),
or attack-emulation events (Mordor, BOTSv3). It is the first benchmark, to our
knowledge, that jointly covers Active Directory identity hygiene state, endpoint
patch posture, vulnerability exposure, and telemetry freshness under controlled
conditions with an openly released, reproducible synthetic generator.
